\section{Stakeholders and System Actors}
\label{sec:stakeholders}

The primary stakeholder of this platform is the Reading the Reader research initiative at DTU Compute \cite{dtucompute2025rtr}, which commissioned it to support controlled studies of adaptive typographic interventions on reading behaviour.
The initiative's supervisors served as the requirements authority throughout the project: they participated in the user story elicitation sessions described in Sec.~\ref{sec:user-stories}, reviewed scope and priority, and provided the research constraints that shaped the architecture.
A secondary stakeholder group consists of future research teams and external contributors who may extend the platform with new intervention strategies or decision algorithms; their interests motivate the modularity and documented integration protocol that run through the design.

Five actors interact directly with the system at runtime.

\subsection{Researcher}
\label{sec:stakeholder-researcher}

The researcher is the primary human operator and the stakeholder around whose workflow the platform is principally organised.
In the context of a controlled reading study, the researcher is responsible for the complete experimental lifecycle: preparing reading materials and experiment templates, enrolling participants and recording their demographic attributes, ensuring the sensing hardware is correctly connected and calibrated, and conducting the session from start to finish.

During an active session, the researcher operates from a physically separate screen, observing the participant's reading behaviour through a mirrored live view without being present at the participant's display.
From this interface, the researcher must be able to monitor session health, manually apply typographic interventions, approve or reject proposals submitted by an automated decision strategy, and pause and resume automated decision-making.
Once a session concludes, the researcher must be able to export session data for external analysis and replay recorded sessions to review gaze behaviour and intervention history at a later point.

Because the researcher coordinates every aspect of an experiment, minimising operator cognitive load is a primary concern.
A guided, step-by-step setup workflow with explicit readiness gating and a unified live control panel are therefore core requirements, not peripheral usability enhancements.

\subsection{Participant}
\label{sec:stakeholder-participant}

The participant is the person recruited by the researcher to take part in a reading experiment.
Their role within the system is deliberately narrow: they read Markdown-formatted text presented in the reading interface while their gaze behaviour is recorded.
The participant does not interact with system configuration, session management, or any researcher-facing controls.

During a session, the participant may be exposed to typographic adaptations (changes to font size, line width, line height, letter spacing, font family, colour palette, or theme mode) applied by the researcher or by an automated strategy.
A critical requirement arising from this role is that such adaptations must not disrupt reading flow or cause the participant to lose their place in the text, as perceptible disruption would confound the experimental results.
Following each reading item, the participant may also be asked to answer comprehension quiz questions, the results of which are recorded as part of the session data.

\subsection{External Module Provider}
\label{sec:stakeholder-provider}

The external module provider is a third-party process that connects to the platform to supply decision logic or eye-movement analysis without requiring modifications to the core system.

A central research motivation for this platform is the ability to evaluate different decision and analysis strategies under comparable experimental conditions.
This actor motivates the need for a stable and documented integration protocol so that different strategies (rule-based, machine-learning-based, or otherwise) can be substituted without changes to the core system.
The platform must support a range of execution modes, from fully autonomous application of proposals to researcher-gated approval, so that the degree of automation can itself be varied as an experimental condition.
The system must also continue operating normally if a provider disconnects mid-session, ensuring that a provider failure does not abort an ongoing experiment.

\subsection{Eye Tracker Device}
\label{sec:stakeholder-eyetracker}

The eye tracker device is the hardware sensing actor that supplies the raw gaze stream on which the platform's adaptive behaviour depends.
In the context of this project the device is a Tobii screen-based eye tracker \cite{tobii2025fusion}, though the system must support hardware sensing sources through a stable interface so that alternative devices can be substituted without changes to the application layer.
The device streams gaze coordinates, per-eye validity signals, and timestamps at high frequency; it does not initiate any use cases of its own but must be detected, licensed, and calibrated before a session can begin.
A key requirement arising from this actor is that the system must handle device disconnection and reconnection gracefully during a session so that a transient hardware failure does not abort an ongoing experiment.

\subsection{Data Consumer}
\label{sec:stakeholder-consumer}

The data consumer represents the downstream actor (an analyst, a collaborator, or an automated analysis script) that processes the records produced at the end of a session.
This actor does not interact with the running system; their interest is in the structure, completeness, and reproducibility of the exported artefacts.
The requirement that follows is that every session export must be self-contained: it must bundle raw gaze samples, derived events, all intervention records with source attribution and timestamps, comprehension results, and calibration metadata in a machine-readable format so that the full experimental sequence can be reconstructed without access to the live application.
The data consumer also motivates the replay interface, which allows the same exported record to be loaded back into the platform for visual review.

